% Source: https://artofproblemsolving.com/community/c3h2605876_geometry_4
% Stable identity: aops-topic-2605876-q01
Let $f(P)=PA+PB+PC$.

\textbf{Step 1: on a segment, $f$ is largest at an endpoint.} Let $X$, $Y$ be points and $P=(1-t)X+tY$ with $0\le t\le1$. For any point $Q$ we have $Q-P=(1-t)(Q-X)+t(Q-Y)$, so by the triangle inequality
\[QP\le(1-t)\,QX+t\,QY.\]
For $0<t<1$, equality needs $Q-X$ and $Q-Y$ to point the same way (or one of them to be zero), i.e. $Q$ on line $XY$ but not strictly between $X$ and $Y$. Adding the inequalities for $Q=A,B,C$ gives
\[f(P)\le(1-t)f(X)+tf(Y)\le\max\left(f(X),f(Y)\right),\]
and for $0<t<1$ the first inequality is strict, because $A$, $B$, $C$ do not all lie on one line.

\textbf{Step 2: reduce to the vertices.} Let $P\ne A$ be a point of the triangle, and let ray $AP$ meet side $BC$ at $Q$, so that $P$ lies on segment $AQ$ and $Q$ on segment $BC$. By Step 1,
\[f(P)\le\max\left(f(A),f(Q)\right)\qquad\text{and}\qquad f(Q)\le\max\left(f(B),f(C)\right),\]
so $f(P)\le\max\left(f(A),f(B),f(C)\right)$. Moreover the inequality is strict unless $P$ is a vertex: if $P\ne Q$ the first step is strict, and if $P=Q$ is not $B$ or $C$ the second step is strict.

\textbf{Step 3.} At the vertices, $f(A)=b+c$, $f(B)=c+a$ and $f(C)=a+b$. The largest of these omits the shortest side.

Hence $PA+PB+PC$ is largest exactly when $P$ is the vertex opposite the shortest side (if two sides tie for shortest, either of the two opposite vertices; for an equilateral triangle, any vertex), and the maximum is
\[\boxed{a+b+c-\min(a,b,c)},\]
the sum of the two longer sides.

\begin{asy}
usepackage("amsmath");
size(4.5cm);
defaultpen(fontsize(10pt));
pair B=(0,0),C=(2.2,0),A=(1.5,3.0),P=(1.25,1.15);
pair Q=extension(A,P,B,C);
draw(A--B--C--cycle,linewidth(0.8));
draw(P--A,dashed);draw(P--B,dashed);draw(P--C,dashed);draw(P--Q,dotted);
dot(P);dot(Q);
label("$A$",A,N);label("$B$",B,SW);label("$C$",C,SE);label("$P$",P,W);label("$Q$",Q,S);
label("$a$",(1.75,0),S,fontsize(9pt));label("$b$",(A+C)/2,NE,fontsize(9pt));label("$c$",(A+B)/2,NW,fontsize(9pt));
\end{asy}

In the sketch $BC$ is the shortest side, so the maximum is at $A$; $Q$ is where ray $AP$ meets $BC$.
